\SourceLeaf{601}\Needspace{22\baselineskip}\vspace{6mm}
\pdfbookmark[1]{Results of the Observations of the Magnetic Association in 1836.}{resultate-1836-anzeige}
\begin{center}\small Gottingische gelehrte Anzeigen.\quad 1837.\quad June 29.\end{center}

\textit{Results of the observations of the Magnetic Association in the year \textup{1836.} Edited by \textup{\textsc{Carl Friedrich Gauss}} and \textup{\textsc{Wilhelm Weber}}. Göttingen \textup{1838.} Published by Dieterich’s bookshop.} (124 pages in 8vo, with 10 lithographed plates.)

The title of this work identifies only part of its contents, namely the part that gave rise to the work in the first place. These pages have already mentioned several times the association that has formed over the past several years to observe some of the most interesting phenomena of terrestrial magnetism simultaneously at specified times agreed in advance (1834.\ Aug.\ 9. and 1835.\ March 7.\ [pp.\ 518, 529 of this volume]). There is therefore no need to repeat what is known. A large number of places within and outside Germany already take part in this association, and the number continues to increase. The ever more abundant results could no longer be communicated solely by private correspondence. Publication in print had become a necessity for those taking part directly. It was also needed so the results could become the common property of all who take an interest in the natural sciences. At the same time, publication offers the most suitable opportunity to discuss subjects gradually in separate shorter and longer papers issued with it. These include matters directly related to the immediate subject, or indirectly related, such as the instruments, their adjustment and use. as well as other matters that may contribute in any way to a more rigorous scientific foundation for the theory of magnetism and galvanism.

The present first instalment contains graphical representations of the observations of magnetic variation at six scheduled times on just as many lithographed plates, with 46 curves in all from fourteen different observing sites, and the numerical observations themselves for the last three scheduled times. Apart from a historical introduction, however, the larger part of the work consists of the
\SourceLeaf{602}following papers. I. Remarks on the design of magnetic observatories and a description of the instruments to be installed in them. Many readers will be pleased to find that the costs and prices are given too. II. The procedure to be used at the scheduled observation times. III. Extract from three years of observations of magnetic declination at Göttingen. IV. Description of a small apparatus for travellers to measure the intensity of terrestrial magnetism in absolute measure. V. Explanations of the plots (provided here) and numerical readings for the scheduled observations.

The other four lithographed plates give a site plan of the Göttingen magnetic observatory, a perspective view of the observation hall and the instruments installed in it, its floor plan, and accurate illustrations of all the individual parts of the magnetometer.

